\documentclass[11pt,a4paper]{article}

% ---------------------------------------------------------------------------
% Code Architecture & Module Reference
% Copper Alloy Digital Twin -- uncertainty-aware surrogate pipeline
% ---------------------------------------------------------------------------

\usepackage[utf8]{inputenc}
\usepackage[T1]{fontenc}
\usepackage{lmodern}
\usepackage[margin=1in]{geometry}
\usepackage{microtype}
\usepackage{xcolor}
\usepackage{graphicx}
\usepackage{booktabs}
\usepackage{longtable}
\usepackage{enumitem}
\usepackage{titlesec}
\usepackage{amsmath}
\usepackage[colorlinks=true,linkcolor=blue!55!black,urlcolor=blue!55!black,citecolor=blue!55!black]{hyperref}

% --- listings for code ---
\usepackage{listings}
\definecolor{codebg}{rgb}{0.97,0.97,0.97}
\definecolor{codekw}{rgb}{0.10,0.30,0.65}
\definecolor{codecomment}{rgb}{0.40,0.55,0.40}
\definecolor{codestr}{rgb}{0.60,0.30,0.20}
\lstdefinestyle{py}{
  language=Python,
  backgroundcolor=\color{codebg},
  basicstyle=\ttfamily\small,
  keywordstyle=\color{codekw}\bfseries,
  commentstyle=\color{codecomment}\itshape,
  stringstyle=\color{codestr},
  numbers=none,
  breaklines=true,
  frame=single,
  rulecolor=\color{gray!40},
  showstringspaces=false,
  columns=fullflexible,
  keepspaces=true,
}
\lstset{style=py}
\newcommand{\code}[1]{\texttt{\small #1}}
\newcommand{\file}[1]{\texttt{\small #1}}

\titleformat{\section}{\Large\bfseries}{\thesection}{0.6em}{}
\titleformat{\subsection}{\large\bfseries}{\thesubsection}{0.6em}{}

\title{\textbf{Copper Alloy Digital Twin}\\[0.3em]
\large Code Architecture \& Module Reference}
\author{Uncertainty-aware surrogate pipeline}
\date{\today}

\begin{document}
\maketitle

\begin{abstract}
\noindent
This document is a developer-facing map of the codebase that implements the
uncertainty-aware surrogate pipeline and the probabilistic model-based
controller (MBC) for the copper-alloy digital twin. It explains the directory
layout, the responsibility of each module, the public API of the core classes,
and how the pieces compose into the two end-to-end workflows (training a
surrogate; running the controller). It complements --- and does not replace ---
the mathematical specification in \file{copper\_digital\_twin\_v4} and the
inline NumPy-style docstrings. The intended reader is a future maintainer or an
ML engineer joining the project who needs to understand \emph{where} things
live and \emph{why}, not to re-derive the math. A short appendix flags
redundant and deprecated files that are safe to remove.
\end{abstract}

\tableofcontents
\vspace{1em}

% ===========================================================================
\section{Orientation}
% ===========================================================================

The project turns each of nine industrial surrogates (three material properties
$\times$ three process steps) into a model that emits a calibrated Gaussian
predictive distribution rather than a point prediction, and feeds those
distributions to a controller that selects process parameters by
\emph{probability of success} rather than binary feasibility.

Two workflows organise the code:

\begin{description}[leftmargin=1.4em,style=nextline]
  \item[Workflow A --- Train a surrogate.]
    Fit Model A (mean + epistemic) and Model B (aleatoric) with AutoGluon,
    recover ensemble weights, build the predictive distribution, calibrate it,
    and persist a self-contained \emph{bundle}. Driven interactively by the
    \file{annealing\_*\_fixed} notebooks, or in batch by \code{ExperimentRunner}.
  \item[Workflow B --- Run the controller.]
    Load surrogate bundles, sweep a dense \mbox{(temperature, time)} grid,
    predict every cell's distribution, optionally propagate upstream
    uncertainty through a chained surrogate, filter cells by a chance
    constraint, rank by cost, and export the chosen operating points. Driven by
    the \file{mbc\_annealing\_*} notebooks and the standalone scripts.
\end{description}

Throughout, a single rule keeps the math in one place: \textbf{all statistical
logic lives in \code{Modeling}}; every other module orchestrates calls to it.

% ===========================================================================
\section{Repository layout}
% ===========================================================================

\begin{lstlisting}[language=bash,basicstyle=\ttfamily\footnotesize]
.
|-- config/
|   `-- Variables.py            # YAML-backed config accessor (varv.PATHS...)
|-- data/raw/                   # one CSV per surrogate
|-- models/                     # persisted AutoGluon bundles + experiment logs
|-- notebooks/
|   |-- modeling/               # reference per-surrogate pipelines (Workflow A)
|   |   |-- annealing_iacs_fixed.ipynb
|   |   |-- annealing_uts_fixed.ipynb
|   |   `-- cold_drawing_uts_*.ipynb
|   |-- experiments/            # grid sweeps + MBC (Workflow A/B)
|   |   |-- run_experiments_iacs.ipynb
|   |   |-- run_experiments_uts.ipynb
|   |   `-- mbc_annealing_uncertainty.ipynb
|   |-- exploration/            # EDA
|   `-- deprecated/             # superseded notebooks (see Appendix)
|-- scripts/
|   `-- script_mbc_annealing_uncertainty.py   # stakeholder-runnable controller
`-- src/
    |-- modeling/
    |   |-- Modeling.py         # *** core: all UA statistics ***
    |   |-- Evaluation.py       # regression + calibration metrics, rollouts
    |   |-- Experiments.py      # config-driven training runner + logging
    |   `-- MBCInference.py     # bundle loading, grid, MC propagation
    |-- metrics/
    |   `-- ProbabilisticEvaluation.py  # chance constraints, cost, selection
    |-- processing/Processing.py        # dataframe cleaning, validation load
    |-- feature_engineering/
    |   |-- FeatureEngineering.py       # labelling, folds, augmentation
    |   `-- Scaler.py                   # EditMinMaxScaler (deterministic only)
    |-- visualization/Plots.py          # plotting helpers
    |-- DataLoader.py                   # LoaderHelper: pickle/model loading
    `-- DataDumper.py                   # CSV/artefact dumping
\end{lstlisting}

\noindent
A surrogate \emph{bundle} is the unit of persistence and the interface between
the two workflows. It is a directory:

\begin{lstlisting}[language=bash,basicstyle=\ttfamily\footnotesize]
<bundle_dir>/
  artifacts.pkl   # weights, base_model_names, recalibration_c,
                  # variance_floor, ood_ref, y_min, y_max, calibration tables
  model_a/        # AutoGluon TabularPredictor (mean + epistemic)
  model_b/        # AutoGluon TabularPredictor (aleatoric variance, log-space)
\end{lstlisting}

% ===========================================================================
\section{Core module: \code{Modeling}}
% ===========================================================================

\file{src/modeling/Modeling.py} is the heart of the pipeline. Every quantity in
the predictive distribution is computed here, so that notebooks, the experiment
runner, and the controller all share one implementation. The methods group into
five concerns.

\subsection{Model A --- mean and epistemic variance}

\begin{longtable}{@{}p{0.34\textwidth}p{0.60\textwidth}@{}}
\toprule
\textbf{Method} & \textbf{Role} \\
\midrule
\endhead
\code{fit\_model\_a(df, target, features, path, \dots, groups=None, **kwargs)} &
  Train the AutoGluon mean ensemble. \code{groups} names a shared fold-id column
  (passed to the \emph{constructor}, not \code{.fit}) so CV folds can be shared
  with H2O; \code{**kwargs} (e.g.\ \code{sample\_weight}) flow to the
  constructor, and any column-valued kwarg is kept in the training frame. \\
\code{collect\_oof\_base\_learners(predictor)} &
  Stack each base learner's out-of-fold predictions into an $(M,N)$ matrix,
  \emph{excluding} \code{WeightedEnsemble} (so NNLS does not collapse to the
  trivial solution). \\
\code{recover\_ensemble\_weights(oof\_matrix, mu\_oof\_ensemble)} &
  Solve $\min_{w\ge 0}\lVert A w - b\rVert_2$ via NNLS, normalise to
  $\sum w = 1$; returns the weights, a recovery-OK flag, and the max
  reconstruction error. \\
\code{compute\_mu\_and\_epistemic\_variance(preds\_matrix, weights, use\_weights)} &
  Weighted (default) or uniform mean and variance over base learners. The
  uniform branch is the robustness fallback behind \code{use\_weighted\_variance}. \\
\bottomrule
\end{longtable}

\subsection{Model B --- aleatoric variance}

\begin{longtable}{@{}p{0.34\textwidth}p{0.60\textwidth}@{}}
\toprule
\textbf{Method} & \textbf{Role} \\
\midrule
\endhead
\code{build\_aleatoric\_targets(y\_true, mu\_oof, sigma2\_epist\_oof)} &
  Construct $\tilde r^2 = \max\!\big((y-\hat\mu_{\text{oof}})^2 -
  \hat\sigma^2_{\text{epist}},\,0\big)$ and return truncation diagnostics
  (fraction of rows clipped to zero). \\
\code{fit\_model\_b(df, features, r\_tilde\_sq, path, \dots)} &
  Train the variance ensemble on $\log(\tilde r^2 + \epsilon)$ for numerical
  stability. \\
\code{predict\_aleatoric\_variance(predictor\_b, X, variance\_floor)} &
  Back-transform via $\exp$ and apply a floor (a small fraction of
  $\mathrm{Var}(y)$) to avoid degenerate zero variance. \\
\bottomrule
\end{longtable}

\subsection{Full predictive distribution}

\code{predict\_with\_uncertainty(X, predictor\_a, predictor\_b, weights,
model\_names, features, \dots)} is the single entry point that assembles a
prediction. The order of operations is deliberate and matters:

\begin{enumerate}[leftmargin=1.6em,itemsep=0.1em]
  \item per-base-learner predictions $\rightarrow$ weighted $\hat\mu$,
        $\hat\sigma^2_{\text{epist}}$;
  \item \textbf{OOD inflation} (if \code{ood\_ref} given): multiply
        $\hat\sigma^2_{\text{epist}}$ by $m(x)^2$ for off-domain inputs;
  \item add $\hat\sigma^2_{\text{aleat}}$ from Model B;
  \item \textbf{recalibration}: $\sigma \leftarrow c\cdot\sigma$;
  \item \textbf{physical truncation}: reshape to the support
        $[y_{\min},y_{\max}]$, appending \code{mu\_trunc},
        \code{sigma\_trunc}, \code{p\_above\_max}.
\end{enumerate}

\noindent
Returns one row per input with \code{mu}, the three variances,
\code{sigma\_total}, \code{ood\_multiplier}, and (when bounded) the truncated
moments. The convenience pair \code{to\_feature\_frame} /
\code{predict\_target} provides a target-agnostic point prediction that accepts
a dict, Series, or DataFrame and infers features from the predictor.

\subsection{Calibration}

\begin{longtable}{@{}p{0.34\textwidth}p{0.60\textwidth}@{}}
\toprule
\textbf{Method} & \textbf{Role} \\
\midrule
\endhead
\code{empirical\_coverage(mu, sigma, y\_true, alpha)} &
  Fraction of points inside the central $\alpha$-interval
  $[\hat\mu - z_\alpha\sigma,\,\hat\mu + z_\alpha\sigma]$,
  $z_\alpha=\Phi^{-1}(0.5+\alpha/2)$. \\
\code{calibration\_table(mu, sigma, y\_true, alphas)} &
  Coverage and gap (empirical $-$ nominal) at several $\alpha$. \\
\code{fit\_recalibration\_scalar(mu, sigma, y\_true, target\_alpha)} &
  Find scalar $c$ minimising the squared coverage gap at one $\alpha$ (default
  $0.9$). Fitted on the same OOF points it audits --- mildly optimistic, by
  design, in the small-data regime. \\
\bottomrule
\end{longtable}

\subsection{Physical constraints and out-of-domain handling}

\begin{longtable}{@{}p{0.34\textwidth}p{0.60\textwidth}@{}}
\toprule
\textbf{Method} & \textbf{Role} \\
\midrule
\endhead
\code{truncated\_gaussian\_moments(mu, sigma, y\_min, y\_max)} &
  Closed-form mean/std of the Gaussian truncated to the physical support, plus
  \code{p\_above\_max} (mass the untruncated law placed beyond the ceiling --- a
  red-flag diagnostic). \\
\code{pr\_within\_spec\_truncated(mu, sigma, spec\_lo, spec\_hi, \dots)} &
  $\Pr(\text{spec\_lo}\le Y\le\text{spec\_hi})$ under the truncated law,
  renormalised by the truncation mass. This is what the controller consumes for
  feasibility. \\
\code{sample\_truncated(mu, sigma, n\_samples, \dots)} &
  Draw $(K,N)$ Monte Carlo samples from the truncated law; the drop-in
  replacement for \code{rng.normal} in chained propagation, so an impossible
  value never feeds the next surrogate. \\
\code{fit\_ood\_reference(X\_train, features)} &
  Fit a Mahalanobis reference (mean, ridged inverse covariance, $95^{\text{th}}$
  percentile distance $d_{\text{ref}}$). Stored in the bundle. \\
\code{ood\_sigma\_multiplier(X, ood\_ref, gamma, max\_multiplier)} &
  Multiplier $m(x)=\min\!\big(1+\gamma\,\max(0,d(x)-d_{\text{ref}})/d_{\text{ref}},\,
  m_{\max}\big)$ applied to $\sigma_{\text{epist}}$; exactly $1$ inside the
  training cloud. \\
\bottomrule
\end{longtable}

\paragraph{Legacy methods.} \code{scale\_and\_pred},
\code{inverse\_target\_if\_any}, and \code{quick\_linear\_regression} are
retained from the deterministic era for backward compatibility and are not part
of the uncertainty-aware path.

% ===========================================================================
\section{Training orchestration: \code{Experiments}}
% ===========================================================================

\file{src/modeling/Experiments.py} turns a hand-run notebook into a reproducible,
logged sweep. Two classes.

\subsection{\code{ExperimentConfig} (frozen dataclass)}

The immutable, hashable identity of one run. Every knob --- target, features,
presets, fold counts, time limits, calibration $\alpha$, physical bounds, fold
seed --- is a field. Key behaviours:

\begin{itemize}[leftmargin=1.4em,itemsep=0.1em]
  \item \code{config\_hash()} --- deterministic 8-char SHA-1 over all fields
        \emph{except} \code{tag} and \code{base\_tag} (cosmetic grouping), so the
        same configuration never retrains twice.
  \item \code{run\_id} $=$ \code{tag\_\_<hash>} --- the run's folder name.
  \item \code{effective\_base\_tag} --- the parent folder grouping all variants
        of a sweep.
  \item \code{variant(**changes)} --- a copy with fields replaced (lists
        coerced to tuples to stay hashable).
\end{itemize}

\subsection{\code{ExperimentRunner}}

Runs a config end-to-end and maintains the on-disk record. Layout:

\begin{lstlisting}[language=bash,basicstyle=\ttfamily\footnotesize]
models/<process>/experiments/<base_tag>/
  experiments_log.csv      # one row per run (schema-reconciled append)
  list_of_the_best.csv     # top-N runs per metric + run_id -> model_dir map
  <run_id>/
    artifacts.pkl  config.yaml  model_a/  model_b/  leaderboard_autogluon.csv
\end{lstlisting}

\begin{longtable}{@{}p{0.34\textwidth}p{0.60\textwidth}@{}}
\toprule
\textbf{Method} & \textbf{Role} \\
\midrule
\endhead
\code{run\_experiment(df, cfg, df\_val=None, force=False)} &
  Train A and B, recover weights, build aleatoric targets, calibrate, fit the
  OOD reference, optionally score a validation set, persist the bundle, and
  append a log row. Skips (and reloads) if the run already exists unless
  \code{force}. \\
\code{run\_grid(df, base\_cfg, grid, df\_val=None, force=False)} &
  Cartesian product over \code{grid}; each variant is stamped with the base tag
  so they share one folder; failures are logged and the sweep continues. \\
\code{build\_best\_list(cfg, top\_n=5)} &
  Regenerate \file{list\_of\_the\_best.csv}: a column per metric (train + val
  RMSE/MAE/MAPE) listing the best run-ids, plus a \code{run\_id} $\rightarrow$
  \code{model\_dir} lookup for one-step loading. \\
\code{build\_weight\_column} / \code{add\_fold\_column} / \code{dataset\_hash} &
  Data-prep utilities: essay weighting, stratified shared folds (delegated to
  \code{FeatureEngineering}), and a content hash so two runs are known to share
  the same data. \\
\bottomrule
\end{longtable}

\paragraph{Schema-reconciled logging.} \code{\_append\_to\_log} reads the
existing CSV, \code{pd.concat}s the new row (unioning columns, filling absent
ones with \code{NaN}), and rewrites. This prevents the corruption that a raw
append causes when runs differ in column set (e.g.\ with/without validation, or
3- vs 4-feature variants).

% ===========================================================================
\section{The controller: \code{MBCInference} + \code{ProbabilisticEvaluation}}
% ===========================================================================

Workflow B is split into two modules: one that \emph{predicts} distributions
over a grid, and one that turns distributions into \emph{decisions}.

\subsection{\code{MBCInference} (\file{src/modeling/})}

\begin{longtable}{@{}p{0.36\textwidth}p{0.58\textwidth}@{}}
\toprule
\textbf{Method} & \textbf{Role} \\
\midrule
\endhead
\code{load\_surrogate(bundle\_dir)} / \code{load\_surrogates(dict)} &
  Load an \code{artifacts.pkl} + both predictors into a \code{SurrogateBundle}. \\
\code{resolve\_best\_bundle(tag\_dir, metric)} / \code{load\_surrogates\_best} &
  Given a \emph{tag} folder, read \file{list\_of\_the\_best.csv} and pick the
  rank-1 run by a metric (default RMSE) --- so callers never hard-code a run
  hash. \\
\code{build\_param\_grid(operational\_limits, step)} &
  Dense inclusive grid over (temperature, time). \\
\code{build\_state\_grid\_df(material\_properties, param\_grid)} &
  Broadcast the fixed material state across every grid cell. \\
\code{predict\_surrogate(bundle, df\_grid, ood\_gamma)} &
  Run one surrogate's \code{predict\_with\_uncertainty} over the grid, returning
  prefixed columns (\code{iacs\_mu}, \code{iacs\_sigma}, \dots). \\
\code{predict\_all(bundles, df\_grid, ood\_gamma)} &
  Run every surrogate, each using only its own features, merging all
  distribution columns. \\
\code{propagate\_upstream\_to\_downstream(upstream, downstream, df\_grid,
  downstream\_feature, k\_samples=200, \dots)} &
  Single-stage Monte Carlo propagation (spec \S4.2): draw $K$ upstream samples
  from the truncated predictive law, run the downstream surrogate on each, and
  aggregate to a downstream mean and a variance that \emph{includes} the
  propagated upstream uncertainty (exposed as \code{*\_sigma2\_propagated}). \\
\bottomrule
\end{longtable}

\noindent
\code{SurrogateBundle} is a small dataclass holding \code{target},
\code{features}, both predictors, the \code{artifacts} dict, and a
\code{short\_name} property (target stem without \code{\_final}).

\subsection{\code{ProbabilisticEvaluation} (\file{src/metrics/})}

The distribution-aware replacement for the deterministic feasibility and cost
logic.

\begin{longtable}{@{}p{0.36\textwidth}p{0.58\textwidth}@{}}
\toprule
\textbf{Method} & \textbf{Role} \\
\midrule
\endhead
\code{add\_success\_probabilities(df, min\_setpoints, surrogate\_bounds, modl)} &
  Append $\Pr(Y\ge\text{setpoint})$ per target (under the truncated law when
  \code{modl} and bounds are given), plus the combined \code{pr\_success\_all}
  (product over targets). \\
\code{apply\_probabilistic\_setpoints(df, delta, combined\_col)} &
  Keep cells meeting the chance constraint $\Pr \ge 1-\delta$. The probabilistic
  counterpart of the deterministic \code{apply\_min\_setpoints}. \\
\code{add\_cost\_function} / \code{add\_cost\_function\_physics} &
  Normalised operational cost (lower $T$, lower $t$) and a thermal-budget proxy
  ($T\!\cdot\!t$). \\
\code{add\_risk\_adjusted\_value(df, short\_name, kappa)} &
  Lower confidence bound $\hat\mu-\kappa\hat\sigma$ for risk-aware ranking
  (prefers tight-$\sigma$ cells among equal means). \\
\code{select\_lowest\_by\_features} / \code{select\_highest\_by\_features} &
  Pick the best row per criterion (lowest cost/time; highest probability). \\
\code{filter\_by\_temperature\_time\_pairs} &
  Targeted inspection of explicit operating points. \\
\bottomrule
\end{longtable}

\paragraph{$\delta$ versus $\kappa$.} They are deliberately decoupled. $\delta$
is the \emph{filter}: a cell is admissible iff $\Pr(\text{success})\ge 1-\delta$.
$\kappa$ is the \emph{ranking knob}: among admissible cells it discounts the
mean by $\kappa$ standard deviations to prefer robust guarantees. $\delta$ is
binary (pass/fail); $\kappa$ only reorders. A common, parameter-free alternative
to a $\kappa$-LCB is to rank directly by \code{pr\_success\_all}, which already
folds $\mu$ and $\sigma$ together in a calibrated way.

% ===========================================================================
\section{Supporting modules}
% ===========================================================================

\begin{description}[leftmargin=1.2em,style=nextline]
  \item[\code{Evaluation} (\file{src/modeling/}).]
    Target-agnostic metrics: \code{one\_step\_metrics} (RMSE, MAE, MAPE, $R^2$,
    optional coverage), \code{metrics\_on\_dataframe} (the same taking a
    validation DataFrame + target name), and the cold-drawing rollout family
    (\code{rollout\_single\_group}, \code{rollout\_all\_groups},
    \code{final\_state\_metrics}, \code{pass\_by\_pass\_curve},
    \code{compounding\_ratio}) for recursive multi-pass evaluation.
  \item[\code{Processing} (\file{src/processing/}).]
    \code{df\_to\_float} (clean/cast with precision-tag stripping),
    \code{load\_validation\_data\_chimie\_paris} (the held-out validation
    loader), \code{reorder\_columns}, and \code{filter\_cold\_drawing\_rules}
    (monotonicity constraints between pre/post properties).
  \item[\code{FeatureEngineering} (\file{src/feature\_engineering/}).]
    Element labelling (\code{label\_element}), ratio masks
    (\code{add\_ratio\_mask\_column}), the shared stratified-fold builder
    (\code{add\_stratified\_fold\_column}, used by both notebooks and
    \code{ExperimentRunner}), and the simulation-expansion / outlier utilities
    used in data preparation.
  \item[\code{Scaler} (\file{src/feature\_engineering/}).]
    \code{EditMinMaxScaler} (fixed-bound min--max with optional $\log1p$) and
    \code{ScalerHelpers}. Used by the \emph{deterministic} path only; AutoGluon
    preprocesses internally, so the uncertainty-aware pipeline does not scale
    inputs.
  \item[\code{Variables} (\file{config/}).]
    Loads \file{variables.yaml} and exposes it as attributes
    (\code{varv.PATHS.models}, etc.), self-refreshing on access.
  \item[\code{DataLoader} / \code{DataDumper} / \code{Plots}.]
    \code{LoaderHelper} loads pickles and models; \code{DataDumper} writes CSVs
    and selections; \code{Plots} holds plotting helpers used in notebooks.
\end{description}

% ===========================================================================
\section{End-to-end: how a prediction is made}
% ===========================================================================

The minimal client (from the README) ties the core pieces together:

\begin{lstlisting}
from src.modeling.Modeling import Modeling
from autogluon.tabular import TabularPredictor
import pickle, pandas as pd

modl = Modeling()
predictor_a = TabularPredictor.load("models/annealing_iacs/<run>/model_a")
predictor_b = TabularPredictor.load("models/annealing_iacs/<run>/model_b")
art = pickle.load(open("models/annealing_iacs/<run>/artifacts.pkl", "rb"))

preds = modl.predict_with_uncertainty(
    X=pd.DataFrame({"purity":[99.95], "iacs":[98.0],
                    "temperature":[573], "time":[60]}),
    predictor_a=predictor_a, predictor_b=predictor_b,
    weights=art["weights"], model_names=art["base_model_names"],
    features=art["features"], variance_floor=art["variance_floor"],
    recalibration_c=art["recalibration_c"],
    use_weights=art["use_weighted_variance"],
    y_min=art.get("y_min"), y_max=art.get("y_max"),
    ood_ref=art.get("ood_ref"),
)
\end{lstlisting}

\noindent
The controller (Workflow B) does the same per grid cell via
\code{MBCInference.predict\_all}, then layers feasibility, cost, and selection
from \code{ProbabilisticEvaluation}. The standalone
\file{scripts/script\_mbc\_annealing\_uncertainty.py} packages this for
stakeholders behind a single YAML config.

% ===========================================================================
\appendix
\section{Appendix: redundant \& deprecated files}
% ===========================================================================

The following were found during a static read of the codebase (imports traced
across \file{notebooks/}, \file{scripts/}, and \file{src/}). None are imported by
any active notebook or script; removing them is safe and would reduce confusion
for new readers. Verify against your latest working tree before deleting.

\begin{longtable}{@{}p{0.40\textwidth}p{0.54\textwidth}@{}}
\toprule
\textbf{File} & \textbf{Finding / recommendation} \\
\midrule
\endhead
\file{src/metrics/Evaluation.py} &
  \textbf{Duplicate} of \file{src/modeling/Evaluation.py} (same class, methods
  reordered). All active notebooks import \code{from src.modeling.Evaluation};
  none import the \file{metrics} copy. \emph{Delete the \file{metrics} copy} and
  keep a single source. \\
\file{src/modeling/Modeling\_old.py} &
  Superseded deterministic \code{Modeling}. Not imported anywhere.
  \emph{Remove} (history is in git). \\
\file{src/modeling/UncertaintyEstimatorAGL.py} &
  Earlier UA approach (separate \code{sigma\_oof} model on absolute residuals,
  uses scalers). Incompatible with the current Model A/B + NNLS pipeline; not
  imported. \emph{Remove or archive.} \\
\file{src/modeling/UncertaintyEstimatorH2O.py} &
  H2O-era variant of the above; superseded by the AutoGluon migration. Not
  imported. \emph{Remove or archive.} \\
\file{src/processing/ColdDrawingProcessingHelper.py} &
  Not imported by any active notebook/script. Confirm it is not pending for the
  cold-drawing work before removing; otherwise move under a clearly-WIP path. \\
\file{src/metrics/Evaluations.py} (plural) &
  Deterministic-controller metrics (\code{gaussian\_nll},
  \code{select\_models}, \dots), imported only by deprecated/deterministic
  notebooks. Keep only if the deterministic controller is still referenced;
  otherwise archive. \\
\file{notebooks/deprecated/*} &
  \file{OLD\_annealing\_iacs\_benchmark}, \file{annealing\_iacs} (old),
  \file{annealing\_tests}, \file{annealing\_uts\_old},
  \file{annealing\_uts\_fixed\_v1}. Superseded by
  \file{notebooks/modeling/annealing\_*\_fixed.ipynb}. Already isolated under
  \file{deprecated/}; safe to prune once nothing references them. \\
\bottomrule
\end{longtable}

\paragraph{Naming nits (non-blocking).} Two near-homonym pairs invite mistakes:
\file{Evaluation.py} vs \file{Evaluations.py}, and \file{Plots.py} vs
\file{Visualization.py}. Consider consolidating each pair, or renaming the
deterministic-era file with a \code{\_deterministic} suffix, so a new reader can
tell at a glance which belongs to the uncertainty-aware path.

\paragraph{Consistency note.} The legacy \code{scale\_and\_pred} /
\code{quick\_linear\_regression} methods still live inside the
uncertainty-aware \code{Modeling}. They are harmless but slightly blur the
class's single responsibility; a future cleanup could move them to a
\code{LegacyModeling} mixin if the deterministic path is fully retired.

\end{document}